(@@M0@@ \uline{and} @@M1@@ \uline{being assigned weight} @@M2@@), \uline{and these} @@M3@@ \uline{relations generate the ideal of relations among the} @@M4@@.
\end{enumerate}

\sgaStableTarget{sga6.e06.cor.4-4}\textbf{Corollary 4.4.} \uline{The elements of} @@M5@@ \uline{of the form}:
@@M6@@
\uline{form a basis of} @@M7@@ \uline{considered as a module over} @@M8@@.

For the proof of the lemma and the corollary, we refer to [2], Exposé 4, pages 4.19, from which they are borrowed.

\sgaStableTarget{sga6.e06.sec.4.5}4.5.
Let @@M9@@ be a quasi-compact scheme, @@M10@@ an @@M11@@-Module locally free of rank @@M12@@, and @@M13@@ a sequence of integers @@M14@@ with sum @@M15@@. Put @@M16@@; @@M17@@ is endowed with a canonical sequence of quotients @@M18@@; put @@M19@@, @@M20@@, where @@M21@@ is the structural morphism @@M22@@. For @@M23@@, denote by @@M24@@ the kernel of @@M25@@. Let @@M26@@ be the Grothendieck group of locally free sheaves of finite type on @@M27@@ (resp. of perfect complexes). Denote by @@M28@@ the class of @@M29@@, and by @@M30@@ the class of @@M31@@. In the case where @@M32@@ is defined in terms of perfect complexes, the series @@M33@@, @@M34@@ still have an evident meaning, where @@M35@@ is the class of @@M36@@. Finally, note that:
@@M37@@

\sgaStableTarget{sga6.e06.thm.4-6}\textbf{Theorem 4.6.} \uline{With the hypotheses and notation of} 4.5, \uline{the} @@M38@@-\uline{algebra} @@M39@@ \uline{is generated by the} @@M40@@; \uline{the ideal of relations among the} @@M41@@ \uline{is generated by the relations deduced from}:
@@M42@@

